\documentclass[10pt,a4paper]{amsart}
\usepackage[margin=19mm]{geometry}
\usepackage{amsmath,amssymb,mathtools}
\usepackage[scr=rsfs,cal=euler]{mathalfa}
\usepackage{xfrac}
\usepackage[unicode,hidelinks,bookmarks=false]{hyperref}
\newcommand{\ie}{i.e.\ }
\newcommand{\cf}{cf.\ }
\newcommand{\resp}{respectively\ }
\newcommand{\Subsection}{\subsection}
\providecommand{\nobreakdash}{\nobreak\hbox{-}}
\setlength{\emergencystretch}{2em}
\multlinegap0pt
\usepackage{amssymb}
\usepackage{mathtools}
\usepackage[shortlabels]{enumitem}
\usepackage{tikz}
\newtheorem{lemma}{Lemma}
\newtheorem{proposition}{Proposition}
\usepackage{cleveref}
\crefname{lemma}{Lemma}{Lemmas}
\crefname{proposition}{Proposition}{Propositions}
\title{Undecidability of Finite Orbit Recognition in Polynomial Maps}
\author{Gwangyong Gwon}
\date{Source: arXiv:2510.00412v3 (CC BY 4.0)}
\begin{document}
\maketitle

\noindent\emph{Source and licence.} Undecidability of Finite Orbit Recognition in Polynomial Maps. Version arXiv:2510.00412v3 (CC BY 4.0), distributed by arXiv under the Creative Commons Attribution 4.0 International licence, http://creativecommons.org/licenses/by/4.0/, as recorded in the arXiv licence metadata for this exact version and shown on the abstract page https://arxiv.org/abs/2510.00412v3. Full licence notice: \texttt{licenses/arxiv-2510.00412-CC-BY-4.0.txt}.\par\smallskip

\noindent\textit{Editorial scope.} Counted unit: the article's proposition prop\_ofrt (source0.tex lines 318--323). The article's complete original proof of it is delivered with it (source0.tex lines 325--331, one \texttt{proof} environment of 7 lines). No mathematics is written, repaired or re-derived: the statement and the proof are the article's own lines, in the article's own notation, and no retained line is substituted or re-flowed.  Retained uncounted as the source-local context the proof needs: lines 254--256; lines 303--308.  Nothing was omitted from any retained span, so the package carries no omission record.  No retained line contains a citation, so the wrapper carries no bibliography and no bibliography entry.  No retained line contains a figure environment, an image-inclusion command or drawing code, so no image is delivered and the package declares no asset dependency.  Editorial formatting only, in addition: an AMS \texttt{amsart} preamble that restates the article's own declarations and macros used by the retained lines, namely the environments \texttt{lemma}, \texttt{proposition} on separate counters, together with the standard packages those lines need.  The retained environments print in this document as Lemma 1, Proposition 1 in order of appearance, whereas the article prints the same items with the numbering of its own full document.  The complete native source of the article as distributed in the arXiv e-print for this exact version is bundled unchanged as \texttt{sources/186134/source0.tex}.
\begin{lemma}[Injectivity]\label{lem_inj}
The encoding $E$ is injective on valid Turing configurations.
\end{lemma}

\begin{lemma}[Simulation lemma]\label{lem_sim}
For every valid Turing configuration $C$,
\[F_N(E(C))=E(T_N(C)).\]
Consequently, for every $t\ge0$,
\[F^t_N(E(C_0))=E(T^t_N(C_0)).\]
\end{lemma}

\begin{proposition}\label{prop_ofrt}
Let $N$ be a deterministic one-tape Turing machine and let $C_0$ be a valid finite-support initial configuration.
Then
\[\mathcal{O}_{F_N}(E(C_0))\text{ is finite}\longleftrightarrow\{T^t_N(C_0):t\ge0\}\text{ is finite.}\]
Equivalently, the polynomial orbit is finite if and only if the Turing trajectory is eventually periodic.
\end{proposition}

\begin{proof}
By \Cref{lem_sim}, every point in the polynomial orbit is the encoding of the corresponding Turing configuration.
By \Cref{lem_inj}, equality of encoded points is equivalent to equality of Turing configurations.
Thus
\[F^k_N(E(C_0))=F^l_N(E(C_0))\quad\longleftrightarrow\quad T^k_N(C_0)=T^l_N(C_0).\]
For a deterministic self-map, the forward orbit is finite if and only if some pair of iterates repeats, which is equivalent to eventual periodicity.
\end{proof}

\end{document}
